% ============================================================
% 第 22 章 直流电机的磁场（教材第六篇，印刷页 244--253）
% 节名：第一节 直流电机的励磁方式；第二节 空载时直流电机的磁场；
%       第三节 负载时直流电机的磁场及电枢反应；第四节 直流电机的感应电动势和电磁转矩
% ============================================================
\chapter{直流电机的磁场}

直流电机的磁场由\textbf{主磁极的励磁绕组}建立（这部分称为主磁场），负载后电枢绕组电流又产生\textbf{电枢磁场}，两者叠加即为气隙磁场。电枢磁场对主磁场的影响称为\textbf{电枢反应}，它是直流电机一系列运行特性（转矩、转速、换向）的根源。

\section{直流电机的励磁方式}

按励磁绕组与电枢绕组的连接方式，直流电机分为四类。

\subsection{他励}
励磁绕组由\textbf{独立的直流电源}供电，与电枢回路无电的联系。特点是励磁电流不受电枢电压影响，便于调节，常用于需要宽范围调速的场合（如轧钢电动机）及发电机试验。

\subsection{并励}
励磁绕组与电枢绕组\textbf{并联}，励磁电流 $I_f=U/R_f$ 取自电枢端电压。并励电动机的励磁电流随端电压变化，具有"自励"特性；并励发电机是最常用的直流发电机。

\subsection{串励}
励磁绕组与电枢绕组\textbf{串联}，励磁电流即为电枢电流 $I_f=I_a$。串励电动机起动转矩大、过载能力强，但轻载时转速极高（"飞车"），绝不能空载运行。

\subsection{复励}
主磁极上有两套励磁绕组，一套与电枢并联（并励绕组），一套与电枢串联（串励绕组）。若两者磁动势方向相同称\textbf{积复励}（常用），相反称\textbf{差复励}。复励兼顾了并励的稳定性和串励的大转矩。

\begin{center}
\begin{tabular}{@{}p{1.8cm}p{3.4cm}p{3.0cm}p{4.4cm}@{}}
\toprule
励磁方式 & 连接 & 励磁电流 & 特点与用途 \\
\midrule
他励 & 独立电源 & $I_f$ 独立可调 & 调速范围宽，轧钢、试验 \\
并励 & 与电枢并联 & $I_f=U/R_f$ & 特性较硬，常用发电机、机床 \\
串励 & 与电枢串联 & $I_f=I_a$ & 起动转矩大，电车、起重机；忌空载 \\
复励 & 并励 + 串励 & 两者合成 & 兼顾稳定性与起动转矩 \\
\bottomrule
\end{tabular}
\end{center}

\section{空载时直流电机的磁场}

空载（电枢电流为零或很小）时，气隙磁场完全由主磁极励磁磁动势建立。

\subsection{主磁场分布}
主磁极的极靴做得比磁极本体宽，使极靴下气隙较均匀、极间气隙很大。因此气隙磁密沿圆周的分布接近\textbf{平顶波}：极靴下磁密基本恒定，极间迅速下降。工程上取\textbf{平均磁密} $B_{av}$ 进行计算，它与最大磁密的关系由\textbf{极弧系数} $\alpha$ 描述：$B_{av}=\alpha B_\delta$（$\alpha$ 一般 0.6～0.75）。

每极磁通
\begin{equation}
\Phi=B_{av}\,\tau\,l=\alpha B_\delta\tau l ,
\end{equation}
其中 $\tau=\pi D_a/(2p)$ 为极距，$l$ 为电枢有效长度。

\subsection{磁化曲线}
空载时改变励磁电流 $I_f$，端电压（或电枢电动势）随之变化，得\textbf{空载特性} $E_0=f(I_f)$。它实质上就是电机的磁化曲线：起始段近似直线（气隙线），随后因铁芯饱和而弯曲。磁路饱和程度直接影响电枢反应的大小与换向的好坏。

\subsection{漏磁}
主磁极的磁通除穿过气隙进入电枢的主磁通外，还有一部分不经电枢、在磁极之间直接闭合，称为\textbf{漏磁通}。漏磁使磁极与机座需要更大的截面，用\textbf{漏磁系数} $\sigma=(\Phi+\Phi_\sigma)/\Phi\approx1.1\sim1.2$ 描述。

\section{负载时直流电机的磁场及电枢反应}

\subsection{电枢磁场}
电枢绕组中流过电流 $I_a$ 时产生电枢磁动势。由于电刷通常放在\textbf{几何中性线}上，电枢磁动势的轴线与主磁极轴线相差 $90^\circ$ 电角度，故其性质为\textbf{交轴电枢反应}。

\subsection{交轴电枢反应的影响}
电枢磁场与主磁场叠加后：
\begin{itemize}
  \item 主磁极的\textbf{前极尖}（进入端）磁通增强、\textbf{后极尖}（退出端）磁通削弱，气隙磁密分布\textbf{畸变}，不再是平顶波；
  \item 由于铁芯饱和，增强的一侧饱和程度加深、增磁有限，而削弱的一侧去磁作用充分，故合成磁通\textbf{略有减小}（一般减小百分之几）；
  \item 磁密的零点（\textbf{物理中性线}）不再与几何中性线重合，而是发生偏移。对\textbf{发电机}，物理中性线\textbf{顺着电枢旋转方向}偏移；对\textbf{电动机}，\textbf{逆着旋转方向}偏移。
\end{itemize}

\subsection{直轴电枢反应}
若电刷不在几何中性线上（例如为改善换向而移刷），电枢磁动势将出现\textbf{直轴分量}：电刷顺转向移动时产生去磁作用，逆转向移动时产生助磁作用。此外，换向极绕组若设计不当也会引入直轴磁动势。

\subsection{电枢反应的不良后果}
\begin{enumerate}
  \item 气隙磁密畸变使\textbf{换向条件恶化}，电刷下产生火花（第 25 章）；
  \item 磁通略有减小，使电动势、转矩下降，需略微增大励磁补偿；
  \item 对大型电机，畸变严重时可能引起环火。
\end{enumerate}

\subsection{抑制措施}
\begin{itemize}
  \item 装\textbf{换向极}：产生直轴磁动势抵消交轴电枢反应，并改善换向（第 25 章）；
  \item 装\textbf{补偿绕组}：嵌在主磁极极靴的槽内，与电枢绕组串联，其磁动势直接抵消交轴电枢反应磁动势，用于负载重、过载大、换向困难的大型电机。
\end{itemize}

\begin{center}
\begin{tabular}{@{}p{2.2cm}p{4.2cm}p{5.0cm}@{}}
\toprule
项目 & 空载 & 负载 \\
\midrule
磁场来源 & 仅励磁磁动势 & 励磁磁动势 + 电枢磁动势 \\
气隙磁密 & 平顶波（极靴下均匀） & 畸变，前后极尖一边增一边减 \\
磁通量 & $\Phi_0$ & $\Phi\approx(0.95\sim0.98)\Phi_0$（略减） \\
物理中性线 & 与几何中性线重合 & 偏移（发电机顺转向、电动机逆转向） \\
后果 & — & 换向恶化、火花、转矩略降 \\
\bottomrule
\end{tabular}
\end{center}

\corepoint{直流电机电枢反应的三句话必须背熟：① 电刷在几何中性线上时是\textbf{交轴电枢反应}，只使磁场\textbf{畸变}，不改变磁通方向；② 畸变的具体表现是"前极尖增磁、后极尖去磁"，因饱和而\textbf{净磁通略减}，\textbf{物理中性线偏移}；③ 偏移方向：发电机\textbf{顺}转向、电动机\textbf{逆}转向。凡是问"电枢反应有什么影响"，按"畸变 $\rightarrow$ 磁通略减 $\rightarrow$ 中性线偏移 $\rightarrow$ 换向变坏"这四步答。}

\section{直流电机的感应电动势和电磁转矩}

\subsection{感应电动势}
电枢旋转时，绕组切割气隙磁通感应电动势。设电枢总导体数为 $N$、并联支路对数为 $a$、每极磁通为 $\Phi$、转速为 $n$，则电枢电动势为
\begin{equation}
E_a=C_e\Phi n,
\end{equation}
其中
\begin{equation}
C_e=\frac{pN}{60a}
\end{equation}
为\textbf{电动势常数}（由电机结构决定）。

\keyformula{E_a=C_e\Phi n,\qquad C_e=\frac{pN}{60a}}

\subsection{电磁转矩}
电枢电流 $I_a$ 在磁场中受力产生转矩：
\begin{equation}
T=C_T\Phi I_a,
\end{equation}
其中
\begin{equation}
C_T=\frac{pN}{2\pi a}
\end{equation}
为\textbf{转矩常数}。两者关系为
\begin{equation}
C_T=\frac{60}{2\pi}C_e=9.55\,C_e .
\end{equation}

\keyformula{T=C_T\Phi I_a,\qquad C_T=\frac{pN}{2\pi a}=9.55C_e}

\subsection{电磁功率}
由 $P_{em}=E_aI_a$ 与 $P_{em}=T\Omega$ 可得
\begin{equation}
E_aI_a=C_e\Phi n\cdot I_a=\frac{C_T\Phi I_a\cdot2\pi n}{60}=T\Omega ,
\end{equation}
即电磁功率既等于电枢吸收（或发出）的电功率，又等于转换为机械（或电）功率的那部分——这是第 2 章 $ei=fv$ 的直流电机版本。

\begin{center}
\begin{tabular}{@{}lll@{}}
\toprule
量 & 表达式 & 说明 \\
\midrule
电枢电动势 & $E_a=C_e\Phi n$，$C_e=pN/(60a)$ & 与转速、磁通成正比 \\
电磁转矩 & $T=C_T\Phi I_a$，$C_T=pN/(2\pi a)$ & 与电流、磁通成正比 \\
常数关系 & $C_T=9.55C_e$ & 由 $\Omega=2\pi n/60$ 导出 \\
电磁功率 & $P_{em}=E_aI_a=T\Omega$ & 机电能量转换的桥梁 \\
\bottomrule
\end{tabular}
\end{center}

\corepoint{$E_a=C_e\Phi n$ 与 $T=C_T\Phi I_a$ 是直流电机的"两把钥匙"，后面所有特性都由这两式推出。注意三点：① $\Phi$ 是\textbf{每极磁通}，不是总磁通；② $C_e$、$C_T$ 只与结构（$p$、$N$、$a$）有关，与运行状态无关；③ 由 $U=E_a+I_aR_a$（电动机）或 $U=E_a-I_aR_a$（发电机）可见，$E_a$ 与 $U$ 之间的差就是电枢电阻压降——这一条是分析机械特性与起动电流的根基。}
